\documentclass[12pt]{article}

%%% Shared preamble (packages, fonts, layout)
\usepackage{linguistics-preamble}
\addbibresource{../../Bibliography/references.bib}

%%% Title
\title{Control and Lowering}
\author{Michael Dover \& Masaya Yoshida}
\date{\today}

\begin{document}
\maketitle

\section{Introduction}\label{sec:intro}

\citet{hornstein1999} proposes that obligatory control is an instance of A-movement. On this view, a sentence like \emph{John wants to leave} is derived by merging \emph{John} as the subject of the embedded infinitival and then raising it to the matrix subject position, exactly as in raising constructions like \emph{John seems to be happy}. The difference between the two reduces to theta-theory: in control, the moved DP receives a theta-role in both its base and landing positions, while in raising, it is theta-marked only in the base position. This \textsc{Movement Theory of Control} (MTC) achieves a significant simplification of the grammar by eliminating PRO as a primitive and unifying control with the independently necessary operation of A-movement.

\citet{neeleman2006} challenge this unification on empirical grounds. They observe that quantifier lowering --- the reconstruction of a higher quantifier into a lower position for purposes of scope --- behaves differently in raising and control structures. Specifically, quantifier lowering in control complements is blocked by any intervening quantifier, whereas in raising structures it is blocked only by negative operators and quantifiers of the form \emph{at most N}. Since the MTC treats both constructions as A-chains, it appears to predict that reconstruction should be subject to identical conditions in both environments. The observed asymmetry, \citeauthor{neeleman2006} argue, shows that the dependency underlying control cannot be the same as the dependency underlying raising.

In this paper, we argue that this conclusion does not follow. The asymmetry that \citeauthor{neeleman2006} identify is real, but it is compatible with --- and indeed expected under --- a version of the MTC that takes seriously the role of phases in the derivation. It is standardly assumed that control infinitivals are CPs, while raising infinitivals are TPs \citep{chomsky2001}. If CP constitutes a phase boundary, then the complement of C is transferred to the interfaces before the matrix derivation proceeds. This creates a locality domain within which reconstruction must be resolved, and it is this domain, we argue, that is responsible for the heightened sensitivity to intervening quantifiers in control structures. In raising complements, where no phase boundary intervenes between the base and surface positions of the moved DP, reconstruction is subject only to the more general constraints on scope interaction in A-chains.

The paper is organized as follows. In \S\ref{sec:nt}, we present \citeauthor{neeleman2006}'s argument in detail, laying out the empirical evidence for quantifier lowering in control and the asymmetries between raising and control that form the basis of their challenge to the MTC. In \S\ref{sec:response}, we develop our phase-based response, showing how the independently motivated distinction between CP and TP complements derives the observed contrasts without abandoning the core insight of the movement analysis.

\section{Neeleman \& Truswell's argument against the Movement Theory of Control}\label{sec:nt}

In this section, we present the argument of \citet{neeleman2006} against a unified movement analysis of raising and obligatory control. Their argument proceeds in two stages. First, they establish that scope ambiguities in control complements are best explained by a rule of quantifier lowering (reconstruction), rather than by quantifier raising out of the infinitival clause. Second, they show that this lowering operation is subject to different intervention conditions in raising and control structures, and conclude that the two cannot instantiate the same syntactic relation.

\subsection{The case for quantifier lowering in control}\label{sec:nt-lowering}

\citeauthor{neeleman2006} begin with the observation that sentences like (\ref{ex:nt-amb}) are scopally ambiguous.
%
\ex\label{ex:nt-amb}
Someone tried [\emph{e} to read every book].\\
$\exists > \forall$;\quad $\forall > \exists$
\xe
%
There are two ways to derive the inverse-scope reading in which the universal takes scope over the existential. On a \emph{quantifier raising} analysis, the universal raises out of the infinitival clause, as in (\ref{ex:nt-qr}). On a \emph{quantifier lowering} analysis, quantifier raising is strictly clause-bound, and the existential is instead lowered (reconstructed) into the subject position of the infinitival, placing it within the scope of the universal. The relevant steps of this derivation are given in (\ref{ex:nt-ql}).
%
\pex\label{ex:nt-two-analyses}
\a\label{ex:nt-qr} {[Every book]$_i$} [someone tried [\emph{e} to read $t_i$]]
\a\label{ex:nt-ql}
\vtop{\halign{#\hfil\cr
\quad (i)\hskip1em [\_\_ tried [someone to read every book]]\cr
\quad (ii)\hskip1em [\_\_ tried {[[every book]$_i$} [someone to read $t_i$]]]\cr
}}
\xe
%
These two analyses make different empirical predictions. The quantifier-raising analysis in (\ref{ex:nt-qr}) is permissive: it allows any quantifier in the embedded clause to take scope over any quantifier in the matrix clause. The quantifier-lowering analysis in (\ref{ex:nt-ql}) is selective: only the category that controls the infinitival subject can be lowered, and therefore only that category can scope below elements in the embedded clause. As \citeauthor{neeleman2006} show, the data support the selective, lowering-based analysis.

Consider the following examples, involving object control in (\ref{ex:nt-objctrl}) and subject control in (\ref{ex:nt-subjctrl}).
%
\pex\label{ex:nt-objctrl}
\a\label{ex:nt-objctrl-a} Mary has persuaded someone [\emph{e} to read every book on the reading list].\\
$\exists > \forall$;\quad $\forall > \exists$
\a\label{ex:nt-objctrl-b} Someone has persuaded Mary [\emph{e} to read every book on the reading list].\\
$\exists > \forall$;\quad $^*\forall > \exists$
\xe
%
\pex\label{ex:nt-subjctrl}
\a\label{ex:nt-subjctrl-a} Someone promised Mary [\emph{e} to read every book on the reading list].\\
$\exists > \forall$;\quad $\forall > \exists$
\a\label{ex:nt-subjctrl-b} Mary promised someone [\emph{e} to read every book on the reading list].\\
$\exists > \forall$;\quad $^*\forall > \exists$
\xe
%
The crucial observation is that the examples are ambiguous only when the quantifier in the matrix clause is the controller of the infinitival subject. In the object control cases in (\ref{ex:nt-objctrl}), the universal can scope over the object in (\ref{ex:nt-objctrl-a}), where the object controls, but not over the subject in (\ref{ex:nt-objctrl-b}). In the subject control cases in (\ref{ex:nt-subjctrl}), the pattern reverses: the universal can scope over the subject in (\ref{ex:nt-subjctrl-a}), where the subject controls, but not over the object in (\ref{ex:nt-subjctrl-b}). This is exactly what the lowering analysis predicts and is unexpected on a long-distance quantifier raising account.

Further support for the lowering analysis comes from the observation that a quantifier in a control complement can never take scope over non-arguments in the matrix clause. The adverbials in (\ref{ex:nt-nonadv}) must take wide scope with respect to the embedded universals.
%
\pex\label{ex:nt-nonadv}
\a John probably wants [\emph{e} to buy every book we're thinking of throwing away].\\
probably $> \forall$;\quad $^*\forall >$ probably
\a When in town, John rarely tries [\emph{e} to visit everyone he knows].\\
rarely $> \forall$;\quad $^*\forall >$ rarely
\xe
%
Finally, \citeauthor{neeleman2006} observe that if the matrix quantifier binds an anaphor, the sentence becomes unambiguous, as in (\ref{ex:nt-anaphor}). This is expected under the lowering analysis: anaphors must be in the scope of their antecedents, which is incompatible with lowering the antecedent into the complement clause.
%
\ex\label{ex:nt-anaphor}
Someone promised himself [\emph{e} to read every book on the reading list].\\
$\exists > \forall$;\quad $^*\forall > \exists$
\xe
%
These facts are all consistent with quantifier lowering and difficult to capture under an analysis that relies on quantifier raising out of the control complement.

\subsection{Asymmetries between raising and control}\label{sec:nt-asymmetries}

Having established the case for quantifier lowering, \citeauthor{neeleman2006} turn to the question of whether lowering in control structures and lowering in raising structures are subject to the same conditions. On the Movement Theory of Control \citep{hornstein1999}, both obligatory control and raising involve A-movement, and the lowering operation observed in control would simply be the familiar reconstruction available in A-chains. This predicts that the conditions on lowering should be identical in the two environments.

\citeauthor{neeleman2006} argue that this prediction is not borne out. They formulate the contrast as follows:
%
\pex\label{ex:nt-generalization}
\a\label{ex:nt-gen-ctrl} In a configuration [ Q$_1$ \ldots{} [ Q$_2$ \ldots{} [ PRO$_1$ \ldots{} ]]], Q$_1$ cannot be lowered.
\a\label{ex:nt-gen-raise} In a configuration [ Q$_1$ \ldots{} [ Q$_2$ \ldots{} [ $t_1$ \ldots{} ]]], where Q$_1$ occupies an A-position, Q$_1$ can be lowered, unless Q$_2 \in \{$negation, \emph{at most} $N\}$.
\xe
%
That is, quantifier lowering in control is blocked by \emph{any} kind of intervening quantifier, while quantifier lowering in raising is blocked only by negative operators and numerals of the form \emph{at most N}.

This asymmetry is illustrated by the following minimal pairs. In (\ref{ex:nt-exactly}), the control example permits only the surface-scope reading when an intervening quantifier of the form \emph{exactly N} is present, while the raising counterpart is ambiguous.
%
\pex\label{ex:nt-exactly}
\a At least one student has promised exactly three lecturers [\emph{e} to read every book].\\
at least 1 $>$ exactly 3 $> \forall$;\quad $^*$exactly 3 $> \forall >$ at least 1
\a At least one student seemed to exactly three lecturers [\emph{t} to be reading every book].\\
at least 1 $>$ exactly 3 $> \forall$;\quad {exactly 3 $> \forall >$ at least 1}
\xe
%
Similarly, existential quantifiers block lowering in control but not in raising.
%
\pex\label{ex:nt-some}
\a At least one student has promised some lecturer [\emph{e} to read every book].\\
at least 1 $> \exists > \forall$;\quad $^*\exists > \forall >$ at least 1
\a At least one student seemed to some lecturer [\emph{t} to be reading every book].\\
at least 1 $> \exists > \forall$;\quad {$\exists > \forall >$ at least 1}
\xe
%
Universal quantifiers pattern the same way: they intervene in control but not in raising.
%
\pex\label{ex:nt-every}
\a At least one student has promised every lecturer [\emph{e} to read every book].\\
at least 1 $> \forall > \forall$;\quad $^*\forall > \forall >$ at least 1
\a At least one student seemed to every lecturer [\emph{t} to be reading every book].\\
at least 1 $> \forall > \forall$;\quad {$\forall > \forall >$ at least 1}
\xe
%
Existential quantifiers introduced by \emph{at least N} likewise intervene in control but not in raising.
%
\pex\label{ex:nt-atleast}
\a At least one student has promised at least three lecturers [\emph{e} to read every book].\\
at least 1 $>$ at least 3 $> \forall$;\quad $^*$at least 3 $> \forall >$ at least 1
\a At least one student seemed to at least three lecturers [\emph{t} to be reading every book].\\
at least 1 $>$ at least 3 $> \forall$;\quad {at least 3 $> \forall >$ at least 1}
\xe
%
Where the two environments \emph{do} converge is with negative quantifiers and \emph{at most N}: these block lowering in both raising and control.
%
\pex\label{ex:nt-neg}
\a At least one student has promised no-one [\emph{e} to read every book].\\
at least 1 $>$ neg $> \forall$;\quad $^*$neg $> \forall >$ at least 1
\a At least one student seemed to no-one [\emph{t} to be reading every book].\\
at least 1 $>$ neg $> \forall$;\quad $^*$neg $> \forall >$ at least 1
\xe
%
\pex\label{ex:nt-atmost}
\a At least one student has promised at most three lecturers [\emph{e} to read every book].\\
at least 1 $>$ at most 3 $> \forall$;\quad $^*$at most 3 $> \forall >$ at least 1
\a At least one student seemed to at most three lecturers [\emph{t} to be reading every book].\\
at least 1 $>$ at most 3 $> \forall$;\quad $^*$at most 3 $> \forall >$ at least 1
\xe

\subsection{The structure of the argument}\label{sec:nt-structure}

The overall logic of \citeauthor{neeleman2006}'s argument can be summarized as follows. If obligatory control is A-movement, as \citet{hornstein1999} proposes, then reconstruction in control should be an instance of the general availability of reconstruction in A-chains. This predicts that quantifier lowering should be subject to the same intervention conditions in both environments. The data in \S\ref{sec:nt-asymmetries} show that this prediction is not met: lowering in control is sensitive to any intervening quantifier, while lowering in raising tolerates all interveners except negative operators and \emph{at most N}. \citeauthor{neeleman2006} conclude that any theory wishing to capture these contrasts must distinguish between the relation underlying control and the relation underlying raising, and that the Movement Theory of Control, which identifies the two, is therefore inadequate.

\section{A phase-based account}\label{sec:response}

The argument in \S\ref{sec:nt} rests on the assumption that if control and raising are both instances of A-movement, then reconstruction should be subject to identical conditions in both environments. In this section, we argue that this assumption is unwarranted. The MTC holds that both control and raising involve A-movement, but it does not hold that the two derivations are structurally identical in all respects. In particular, the two differ in the size of the embedded complement: control infinitivals are standardly taken to be CPs, while raising infinitivals are TPs \citep{chomsky2001}. If CP is a phase, this difference has direct consequences for reconstruction, and we show that the asymmetries documented by \citeauthor{neeleman2006} follow.

\subsection{The structural distinction between raising and control}\label{sec:structures}

Under the MTC, a raising sentence like (\ref{ex:raising}) and a control sentence like (\ref{ex:control}) are both derived by A-movement of the subject DP from the embedded clause to the matrix subject position. However, the embedded complements differ in categorial status.
%
\ex\label{ex:raising}
John seems to be reading every book.
\xe
%
\ex\label{ex:control}
John wants to read every book.
\xe
%
In raising, the complement of the matrix verb is a bare TP. There is no CP layer, and crucially, no phase boundary between the base and surface positions of the moved DP. The structure is given schematically in (\ref{ex:raising-tree}).
%
\ex\label{ex:raising-tree}
\begin{forest}
[TP
  [DP$_1$ [\emph{John},roof]]
  [T$'$
    [T]
    [VP
      [V [\emph{seems}]]
      [TP
        [$\langle$DP$_1\rangle$]
        [T$'$
          [T [\emph{to}]]
          [VP
            [\emph{be reading every book},roof]
          ]
        ]
      ]
    ]
  ]
]
\end{forest}
\xe
%
In control, by contrast, the complement of the matrix verb is a CP. On the standard assumption that CP constitutes a phase \citep{chomsky2001}, the complement of C --- namely TP --- is transferred to the interfaces when the phase is completed. The moved DP must pass through the specifier of CP (the phase edge) on its way to the matrix clause. The resulting structure is given in (\ref{ex:control-tree}).
%
\ex\label{ex:control-tree}
\begin{forest}
[TP
  [DP$_1$ [\emph{John},roof]]
  [T$'$
    [T]
    [VP
      [V [\emph{wants}]]
      [CP,name=cp
        [$\langle$DP$_1\rangle$]
        [C$'$
          [C]
          [TP,name=tp,tikz={\node[draw,dashed,inner sep=4pt,fit to=tree]{};}
            [$\langle$DP$_1\rangle$]
            [T$'$
              [T [\emph{to}]]
              [VP
                [\emph{read every book},roof]
              ]
            ]
          ]
        ]
      ]
    ]
  ]
]
\end{forest}
\xe
%
The dashed box around the embedded TP in (\ref{ex:control-tree}) marks the spell-out domain: the material that is transferred to the interfaces when the CP phase is completed. The copy of DP$_1$ at the phase edge (Spec,CP) remains accessible to the continuing derivation, but the copy inside the transferred TP does not.

This structural difference between raising and control is not a stipulation introduced to handle the scope facts. It is independently motivated by a range of considerations, including the distribution of complementizers, the availability of extraction, and the behavior of tense in the two types of complements. What we show in the remainder of this section is that this independently motivated distinction derives the intervention asymmetries documented by \citeauthor{neeleman2006}.

\subsection{Successive-cyclic reconstruction}\label{sec:cyclic-recon}

Our account rests on the following hypothesis about reconstruction in phase-based derivations:
%
\ex\label{ex:hypothesis}
\textsc{Successive-Cyclic Reconstruction}: Reconstruction of a moved DP to a position within a transferred spell-out domain must proceed through the phase edge. Each intermediate step of reconstruction constitutes an independent scope computation.
\xe
%
This hypothesis parallels the standard assumption that movement must proceed successive-cyclically through phase edges \citep{chomsky2001}. Just as a DP cannot move from a position inside a phase to a position outside it without first passing through the phase edge, so too reconstruction --- the interpretive ``undoing'' of movement --- must retrace the path of the derivation step by step.

The consequences of (\ref{ex:hypothesis}) for raising and control are as follows.

\paragraph{Raising.} In raising structures, the A-chain connecting the surface position of DP$_1$ and its base position is contained within a single spell-out domain. No phase boundary intervenes. Reconstruction of DP$_1$ is therefore a single operation: the lower copy is interpreted directly, and the scope of DP$_1$ is evaluated with respect to the other elements in the clause in a single computation. The only quantifiers that block this operation are those that are independently known to induce scope rigidity in A-chains --- namely negative operators and \emph{at most N} \citep[cf.][]{fox1999}. This is schematized in (\ref{ex:raising-recon}).
%
\ex\label{ex:raising-recon}
{[} DP$_1$ \ldots{} Q$_2$ \ldots{} {[}\textsubscript{TP} $\langle$DP$_1\rangle$ \ldots{} $\forall$ \ldots{} {]} {]}\\[4pt]
$\underbrace{\hspace{18em}}_{\text{single spell-out domain: reconstruct in one step}}$
\xe
%
Since there is only one step, Q$_2$ can only block reconstruction if it is inherently scope-rigid.

\paragraph{Control.} In control structures, the A-chain spans two spell-out domains, separated by the CP phase boundary. DP$_1$ has copies in three positions: (i) its surface position (matrix Spec,TP), (ii) the phase edge (embedded Spec,CP), and (iii) its base position (embedded Spec,\emph{v}P/TP). To reconstruct to position (iii), DP$_1$ must first reconstruct to position (ii) --- the phase edge --- and then from (ii) into the transferred domain. Each step is an independent scope computation:
%
\ex\label{ex:control-recon}
{[} DP$_1$ \ldots{} Q$_2$ \ldots{} {[}\textsubscript{CP} $\langle$DP$_1\rangle$ {[}\textsubscript{TP} $\langle$DP$_1\rangle$ \ldots{} $\forall$ \ldots{} {]} {]} {]}\\[4pt]
$\underbrace{\hspace{8em}}_{\text{Step 1}}$\hspace{4.5em}$\underbrace{\hspace{8em}}_{\text{Step 2}}$
\xe
%
\textsc{Step 1}: DP$_1$ reconstructs from its surface position to the phase edge (Spec,CP). This step crosses Q$_2$, which occupies a position in the matrix clause between DP$_1$ and the phase edge. For reconstruction to proceed, DP$_1$ must be interpreted as taking scope below Q$_2$ at this intermediate position. This is an independent scope computation: Q$_2$ has the opportunity to ``trap'' DP$_1$ in a position above it.

\textsc{Step 2}: DP$_1$ reconstructs from the phase edge into the transferred TP, where it takes scope below the universal quantifier.

Crucially, Step 1 is where the additional intervention effects arise. In raising, this step does not exist --- reconstruction proceeds directly from the surface position to the base position within a single domain, and Q$_2$ can only interfere if it is inherently scope-rigid. In control, the intermediate reconstruction to the phase edge creates an additional point at which \emph{any} intervening quantifier can block the process, because the computation at Step 1 involves a genuine scope interaction between DP$_1$ and Q$_2$.

\subsection{Deriving the pattern}\label{sec:deriving}

We are now in a position to derive the full pattern of intervention effects observed by \citeauthor{neeleman2006}.

\paragraph{Raising: only negation and \emph{at most N} block.} In a raising structure like (\ref{ex:raising-some-r}), reconstruction of \emph{at least one student} to its base position is a single operation within one spell-out domain. The intervening existential \emph{some lecturer} does not induce scope rigidity, so it does not block reconstruction. The universal \emph{every book} can therefore take widest scope.
%
\ex\label{ex:raising-some-r}
At least one student seemed to some lecturer to be reading every book.\\
at least 1 $> \exists > \forall$;\quad $\exists > \forall >$ at least 1
\xe
%
When the intervener is a negative quantifier or \emph{at most N}, however, it independently blocks scope reconstruction in A-chains, and the inverse-scope reading is unavailable:
%
\ex\label{ex:raising-neg-r}
At least one student seemed to no-one to be reading every book.\\
at least 1 $>$ neg $> \forall$;\quad $^*$neg $> \forall >$ at least 1
\xe

\paragraph{Control: any intervener blocks.} In a control structure like (\ref{ex:control-some-c}), reconstruction of \emph{at least one student} must proceed in two steps. At Step 1, the subject reconstructs from matrix Spec,TP to Spec,CP, crossing the intervening existential \emph{some lecturer}. This intermediate reconstruction requires \emph{at least one student} to scope below \emph{some lecturer}. Because this is an independent scope computation at the phase boundary, any scope-taking element in the intervening position --- including existentials, universals, and \emph{exactly N} quantifiers --- can trap the subject above it, blocking further reconstruction.
%
\ex\label{ex:control-some-c}
At least one student has promised some lecturer to read every book.\\
at least 1 $> \exists > \forall$;\quad $^*\exists > \forall >$ at least 1
\xe
%
The same holds for universal and \emph{exactly N} interveners:
%
\pex\label{ex:control-others}
\a At least one student has promised every lecturer to read every book.\\
at least 1 $> \forall > \forall$;\quad $^*\forall > \forall >$ at least 1
\a At least one student has promised exactly three lecturers to read every book.\\
at least 1 $>$ exactly 3 $> \forall$;\quad $^*$exactly 3 $> \forall >$ at least 1
\xe

\paragraph{Convergence with negation and \emph{at most N}.} The account correctly predicts that negative quantifiers and \emph{at most N} block reconstruction in both environments. In raising, they block because they are inherently scope-rigid in A-chains. In control, they block for the same reason at Step 1, and additionally would block at Step 2 if Step 1 were to succeed. The two independent sources of blocking converge, yielding the identical behavior in both structures that \citeauthor{neeleman2006} observe.

\subsection{The logic of the argument}\label{sec:logic}

To summarize, the asymmetry between raising and control that \citeauthor{neeleman2006} document does not require us to posit distinct types of syntactic dependencies. It requires only two independently motivated assumptions: (i) control infinitivals are CPs while raising infinitivals are TPs, and (ii) reconstruction, like movement, proceeds successive-cyclically through phase edges. Together, these assumptions predict that reconstruction in control will be subject to an additional layer of locality constraints --- the requirement to reconstruct through the phase edge --- that reconstruction in raising is not subject to.

The MTC is not committed to the claim that every property of A-chains must be identical in raising and control. It is committed to the claim that the same operation --- A-movement --- underlies both. The present account preserves this commitment while deriving the observed differences from the structural contexts in which A-movement applies. The phase boundary in control complements is an independently needed piece of syntactic architecture; once its consequences for reconstruction are taken into account, the asymmetry becomes not a problem for the MTC but a prediction of it.

\printbibliography

\end{document}
